\documentclass[10pt]{article}
\usepackage[margin=1in]{geometry}
\usepackage{amsmath,amssymb}
\usepackage{xurl}
\setlength{\parindent}{0pt}
\setlength{\parskip}{3pt}
\newcommand{\Tr}{\operatorname{Tr}}
\newcommand{\tr}{\operatorname{tr}}
\newcommand{\sources}[1]{\par{\footnotesize\textsc{Sources.} #1\par}}
\begin{document}

\textbf{(a) What first-order elastic amplitudes determine.}
For prescribed static scalar profiles, choose
\[
 A=m\psi^{-z}+i\psi^{-\alpha}D\psi^{-(z+2-\alpha)},\quad
 L=N\psi^z,\quad \alpha=(z+2)(1-\theta),\quad m>0 .
\]
The Hermitian flat seed obeys $d_p^2=\kappa_p^2I$; set
$\epsilon_p=(m^2+\kappa_p^2)^{1/2}$,
$U_p=(m+id_p)/\epsilon_p$ and $v_p=(I,U_p)^T/\sqrt2$.
At $N=\psi=1$ write $\delta\psi=f$, $\delta N=n$, $q=p'-p$.
Differentiating the full polar
block Hamiltonian and projecting on these band states gives
\[
 v_{p'}^\dagger\delta h\,v_p
 =\sqrt{\epsilon'\epsilon}(n_q+zf_q)
   +\tfrac12(\delta A^\dagger U+U'^\dagger\delta A).
\]
At fixed $\alpha$, $\delta A_z=-zf_q(m+id_{p'})$, so the entire
remaining $z$ dependence is
\[
 M_z=zf_q\!\left(\sqrt{\epsilon'\epsilon}
                   -\frac{\epsilon'+\epsilon}{2}\right)I
     =-\frac z2f_q(\sqrt{\epsilon'}-\sqrt\epsilon)^2I .
\]
It vanishes at $\epsilon'=\epsilon$.  Negative-band projection changes
the overall sign and gives the same conclusion for holes.  Thus the
elastic response determines at most $\alpha$; $(z,\theta)=(1,0)$ and
$(2,\tfrac14)$ both give $\alpha=3$.  The exponent is a chosen input.
The separate off-diagonal block is not the band-projected amplitude.
\sources{\path{<internal>/rederivation_final3_selector_filter.tex},
section 1, full-block projection; independent polar-Sylvester route in
\path{<internal>/redteam_final3_selector_filter.tex}.}

\bigskip
\textbf{(b) Nonnegative constrained static energy.}
On a finite open regulator, retain every interior--collar matrix
element and set $\psi_b=1$, $N_b=N_\infty>0$ on the collar $B$.
For positive commuting rulers and a state in the common excitation
form domain, write
\[
 X_x=\tr_2\sum_s\eta_s(|A_s^\dagger|-|A_s|)_{xx},\quad
 e_x=\psi_x\sum_r\langle Q_{rx}^\dagger Q_{rx}\rangle,\quad
 C_x=\psi_xX_x+e_x .
\]
The $Q_{rx}$ are independent of the varied lapse.  Its affinity defines
the same-Hamiltonian boundary energy
$E_{\partial}=\partial_{N_\infty}\langle H_*\rangle=\sum_{b\in B}C_b$.
The two polar moduli have identical singular values, hence
$\sum_xX_x=0$.  Summing $C_x/\psi_x$ therefore gives
\[
 E_{\partial}=\sum_{x\in I\cup B}\frac{e_x}{\psi_x}
                   -\sum_{i\in I}\frac{C_i}{\psi_i}
 \quad\Longrightarrow\quad
 E_{\partial}\big|_{C_I=0}
       =\sum_{x,r}\langle Q_{rx}^\dagger Q_{rx}\rangle\ge0 .
\]
The comparison lapse $N_x^*=N_\infty/\psi_x$ gives the same result:
its vacuum allocation is constant and its energy differs only by
interior clock constraints.  With a finite same-law exterior limit,
$E_{\rm bdy}=\lim E_{\partial}=\sum_xe_x/\psi_x$; collar sources
remain included.  For allocation $N\psi^z$ and its corresponding
positive clock source, the same sum is $\sum_xe_x/\psi_x^z\ge0$,
independently of the vacuum orientation.
\sources{\path{<internal>/rederivation_final3_vacuum_surface.tex},
section 1; independent trace and comparison-lapse reconstructions in
\path{<internal>/redteam_final3_vacuum_surface.tex},
duty 2.}

\newpage
\textbf{(c) The spacing window belongs to the small-spacing branch.}
Fix $m_s^2=M^2r_s$, with
$r=(r_0,r_0+b_1,r_0+b_2,r_0+b_1+b_2)$,
$r_0,b_1,b_2>0$, and $\eta=(+,-,-,+)$.
For $\mu=aM$, $S(k)=\sum_i\sin^2k_i$ and the full-zone average, the
thermodynamic zero-momentum coefficient at uniform ruler $\psi=P$ is
\[
 I(\mu)=\left\langle\cos^2k_1\sum_s
             \frac{\eta_s}{\sqrt{S(k)+\mu^2r_s}}\right\rangle,\qquad
 k_a(P)=\frac{3I(\mu P^2)}{a^2P^4},\qquad
 q_a=-4+2\mu\frac{I'}I .
\]
The rectangle is the positive double difference of
$(S+\mu^2r)^{-1/2}$, so $I>0$ and $q_a$ is continuous.
Changing variables $p_i=\sin k_i$ gives angularly averaged Jacobian
$1+p^2/6+O(p^4)$.  With
$j=\sum_s\eta_sr_s\log r_s>0$ and $s_2=2b_1b_2$,
the grouped radial integral and its differentiated small-$\mu$
expansion yield
\[
 I(\mu)=\frac{j\mu^2}{\pi^2}
       +\frac{s_2\mu^4}{4\pi^2}\log(1/\mu)+O(\mu^4),
 \qquad
 q_a=\frac{s_2}{j}\mu^2\log(1/\mu)+O(\mu^2)>0 .
\]
At large $\mu$, the uniform expansion
$I=(2\mu)^{-1}\sum_s\eta_sr_s^{-1/2}+O(\mu^{-3})$
has positive leading coefficient and gives $q_a\to-6$.
Continuity forces a finite sign change, without uniqueness; a window
based on the positive defect is a small-spacing-branch statement.
At $r_0=0$, with $0\log0=0$, the coefficient argument still holds
but $q_a\to-4$; it supplies no gapped polar construction.
Zero-width rectangles are excluded because $I=0$.
\sources{\path{<internal>/rederivation_final3_vacuum_surface.tex},
sections 3--4, sine-Jacobian route;
\path{<internal>/redteam_final3_vacuum_surface.tex},
duty 3, independent heat-tail route.}

\bigskip
\textbf{(d) The continuum two-derivative vacuum term at full amplitude.}
Take smooth commuting rulers in a compact positive interval, a strictly
positive mass rectangle, and the grouped spin trace on a compatible
closed Dirac domain.  In one continuum valley write
$A=M+f\,\boldsymbol{\sigma}\cdot\nabla$,
$M=m/\psi$, $f=\psi^{-3}$.
After the common measure change, the two polar squares have scalar
potential difference
$v_1-v_0=-f\Delta f+2|\nabla f|^2$
and Pauli-potential square difference
$|q_1|^2-|q_0|^2=-4fM\nabla M\cdot\nabla f$.
One scalar heat insertion and two Pauli insertions give
\[
 \tr_2\langle x|e^{-tAA^\dagger}-e^{-tA^\dagger A}|x\rangle^{[2]}
 =\frac{2}{(4\pi t)^{3/2}}\,
    \nabla\cdot\!\left[t f^{-2}e^{-tM^2}\nabla f\right].
\]
Common scalar terms cancel; the linear connection-curvature insertion
vanishes by its antisymmetry against the isotropic Gaussian.
Group the masses before the square-root heat integral:
$\int_0^\infty t^{-2}\sum_s\eta_se^{-tm_s^2/\psi^2}dt=J/\psi^2$,
where $J=\sum_s\eta_sm_s^2\log(m_s^2/x_{\rm ref})>0$; the reference
cancels by $\sum_s\eta_sm_s^2=0$.
Since $\psi^4\nabla f=-3\nabla\psi$, the eight sine valleys give
\[
 X^{[2]}=\frac{3J}{\pi^2}\Delta\psi,\qquad
 W_{0,\mathrm{bulk}}^{[2]}
 =-k\int\nabla\psi\cdot\nabla(N\psi),\qquad k=\frac{3J}{\pi^2}.
\]
This is the exact two-derivative coefficient, without an amplitude
expansion; integration by parts retains
$k\oint N\psi\,\partial_n\psi$.
\sources{\path{<internal>/rederivation_final3_vacuum_surface.tex},
sections 2--3, heat insertions and arbitrary-background Sylvester route;
\path{<internal>/redteam_final3_vacuum_surface.tex},
duty 4.}

\newpage
\textbf{(e) The finite-coupling ruler row includes reference transport.}
At finite gapped regulator, let $P[\psi]$ be the free occupied
projector and $\mathcal R_P$ retain the quartic monomials with two
excitation creators and two annihilators.  Hold the physical seed
$V=O(g)$ fixed.  For $k_P=[\delta P,P]$ and its Fock lift
$\widehat k_P$, covariance of this selection gives exactly
\[
 (\delta_P\mathcal R_P)V
   =[\widehat k_P,\mathcal R_PV]
                    -\mathcal R_P[\widehat k_P,V].
\]
The chosen energy-sum denominator in the completion gives
$H_{\rm exc}=H_0+\mathcal R_PV+O(g^2)$ at unit lapse.
Thus the displayed transport contributes at order $g$ to the
fixed-physical-state ruler derivative.  Its sign is unrestricted.
For two particle and two hole modes of energy $E$ and total excitation
number $n$, the fixed mass-bilinear seed has moving-frame expression
$-g[c(n-2)+\sqrt{1-c^2}(J_++J_-)]^2$, where
$c=m/\sqrt{m^2+p^2\psi^{-4}}$ and $0<c^2<1$.
Its bright and dark pair energies are
\[
 E_B=2E e^{-g(2-c^2)/E},\qquad E_D=2E e^{-gc^2/E},\qquad
 \partial_{\log\psi}c^2=4c^2(1-c^2).
\]
Here $J_+=a_1^\dagger h_1^\dagger+a_2^\dagger h_2^\dagger$.
Holding $E,g$ and explicit weights fixed, their reference derivatives
are respectively $+8gc^2(1-c^2)e^{-g(2-c^2)/E}$ and
$-8gc^2(1-c^2)e^{-gc^2/E}$.
Wick subtraction changes only degrees zero and two, which
$\mathcal R_P$ erases.  The projector and unit-lapse dressing are
clock independent, so $\partial_{N_x}H_{\rm exc}
=\psi_x\sum_rQ_{rx}^\dagger Q_{rx}$; the dressed annihilators obey
$F_i\Omega_P=0$, leaving the vacuum energy exactly $W_0$.
\sources{\path{<internal>/rederivation_final3_selector_filter.tex},
sections 2--3, transport identity and exact pair spectrum;
\path{<internal>/redteam_final3_selector_filter.tex},
independent CAR and covariant-filter reconstructions.}

\end{document}
